\section{Implementation Details}
\label{app:impl}

All values below are the ones used for the paper's experiments; the exact launch
commands are in the project repository.

\textbf{Excitation and dataset.}
Excitation overrides the converged policy's command with commands drawn uniformly
over the admissible box $[-1,1]^3$ in $(v_x, v_y, \omega_z)$, resampled every 25
low-level steps; with probability $0.3$ the drawn command is scaled by $0.3$,
oversampling the low-magnitude regime that proportional control visits near the
goal. Collection runs 64 parallel environments (8{,}000 steps per environment for
H1). Constant-command windows of 16 low-level steps ($0.32$\,s at the 50\,Hz
control rate) are sliced at stride 4 within episodes, giving the $\approx$25k Go2
windows of Table~\ref{tab:model}. Terrain-relative height labels are recovered
from the policy's own height-scan observation (mean scan return plus its
$0.5$\,m sensor offset).

\textbf{Response model and training.}
$f_\phi$ is a multilayer perceptron with two hidden layers of width 256 (ReLU,
dropout $0.1$), mapping the full policy observation, with the command slice
replaced by the candidate command, to the six response targets; inputs and
targets are standardized. Training uses AdamW (learning rate $10^{-3}$, weight
decay $10^{-4}$), batch size 1{,}024, 200 epochs, keeping the checkpoint with the
best held-out loss. The held-out set is $20\%$ of episode groups, so no episode
straddles the split; Table~\ref{tab:model} reports this held-out set. The
torque-derived mechanical-power channel is identified with the same recipe by a
five-member ensemble.

\textbf{Candidate set.}
The structured set $\mathcal{C}$ around the direct command comprises: uniform
scalings $\{0.4, 0.6, 0.8, 1.2\}\,c_{\text{goal}}$ and $c_{\text{goal}}$ itself;
yaw-rate offsets of $\{\pm0.2, \pm0.4\}$\,rad/s applied to $c_{\text{goal}}$; and
a polar grid of absolute speeds $\{0.25, 0.5, 0.75, 1.0\}$\,m/s at heading
offsets $\{0, \pm0.3, \pm0.6\}$\,rad around the target bearing with yaw-rate
scalings $\{0, 0.5, 1.0\}$ of the direct yaw command. All candidates are clipped
to the admissible box $[-1,1]^3$ and deduplicated, leaving $\approx$70.

\textbf{Controller constants.}
One high-level decision is taken every $E=4$ low-level steps. The gate scalars
shared by all three robots are: correction gain $\alpha=0.5$ with anneal
distance $d_{\text{anneal}}=0.5$\,m (the correction fades out linearly inside
$0.5$\,m of the target), progress floor $\beta=0.9$, and cost margin
$\epsilon=0.1$; the cost channel $\hat C$ is the effort proxy for the main
configuration and predicted mechanical power for the power configuration of
Table~\ref{tab:objective-audit}. The morphology-dependent posture floors are the
eligibility, rescue, and restore fractions of nominal height with the
height-decline tolerance $\delta$: $(0.85, 0.80, 0.90)$ with $\delta=2$\,cm for
both quadrupeds, $(0.93, 0.90, 0.96)$ with $\delta=0.8$\,cm for H1
(Section~\ref{sec:setup}).

\textbf{VGCC pseudocode.}
Algorithm~\ref{alg:vgcc} states one high-level decision. If no candidate is eligible, or the cost margin fails, the controller executes $c_{\text{goal}}$ exactly.

\begin{algorithm}[t]
\caption{Viability-Gated Command Compensation (one high-level step)}
\label{alg:vgcc}
\begin{algorithmic}[1]
\Require frozen policy $\pi_\theta$, response model $f_\phi$, target $g$, correction gain $\alpha$, anneal distance $d_{\text{anneal}}$, and fixed scalars $(\beta,h_{\rm floor},h_{\rm resc},\delta,\epsilon)$
\State Compute body-frame target $g$, direct command $c_{\text{goal}}$ (Eq.~\ref{eq:cgoal}), and the candidate set $\mathcal{C}$ around $c_{\text{goal}}$
\State Batch-predict $f_\phi(o,c)$ for all $c\in\mathcal{C}$; set $\alpha_t=\alpha\operatorname{clip}(d/d_{\text{anneal}},0,1)$
\If{measured posture is below the operating band}
  \State $c\gets c_{\text{goal}}$ \Comment{restore}
\ElsIf{direct control is predicted to breach the posture floor $h_{\rm resc}$}
  \State $c^\ast\gets\arg\max_{\mathcal{C}:\,\hat G\ge 0}\hat h_{\min}$;\quad $c\gets c_{\text{goal}}+\alpha_t(c^\ast-c_{\text{goal}})$ \Comment{rescue}
\Else
  \State $c^\ast\gets\arg\min_{\mathcal{C}\,\text{eligible (Eq.~\ref{eq:vgc-eligible})}}\hat C$
  \If{$\hat C(c^\ast)<(1-\epsilon)\,\hat C(c_{\text{goal}})$ (Eq.~\ref{eq:vgc-cost})}
    \State $c\gets c_{\text{goal}}+\alpha_t(c^\ast-c_{\text{goal}})$ \Comment{compensate}
  \Else
    \State $c\gets c_{\text{goal}}$ \Comment{exact direct fallback}
  \EndIf
\EndIf
\State Execute $\pi_\theta$ with $c$ for $E$ low-level steps
\end{algorithmic}
\end{algorithm}

\textbf{Sampling MPC baseline.}
The MPC of Section~\ref{sec:results-baselines} does \emph{not} recursively use
the six-output response model $f_\phi$: that model has no next-observation
output and cannot be rolled forward. Instead, MPC uses a separately trained
five-member macro-transition ensemble
\[
F_\eta(o_t,c_t)\mapsto
(\widehat{o_{t+4}-o_t},\widehat{\Delta x},\widehat{\Delta y},
\widehat{\Delta\psi},\widehat P,\widehat h_{\min},
\widehat{\theta}_{\max},\widehat{\omega}_{\max}),
\]
where one macro step is four low-level steps ($0.08$\,s). The model has two
hidden layers of width 512 and is trained on 42{,}054 windows from 207 episode
groups with disjoint train/validation/test groups. Each ensemble member
recursively advances its own predicted full observation residual; the command
slice is overwritten by the next planned command before the following query.
At four macro steps ($0.32$\,s), held-out open-loop errors are $3.3$\,cm
displacement MAE, $5.17$\,J cumulative-work MAE, $8.8$\,mm minimum-height MAE,
and $0.756$ normalized full-state RMSE.

MPC enumerates all $3^4=81$ scale sequences from
$\{0.75,0.90,1.00\}$ and executes the first command of the selected sequence.
After the four-step rollout, a second, separately trained task-value network
predicts remaining work, remaining time, and success under direct-control
continuation from the predicted terminal observation and local target. It is
trained on 4{,}298 samples from 144 mixed-controller episodes; held-out metrics
are work MAE $15.7$\,J ($R^2=0.846$), time MAE $0.362$\,s
($R^2=0.654$), and success AUC $0.819$. MPC minimizes ensemble-upper-bound total
work subject to predicted success, time, posture, and optional stability
constraints, with exact direct-sequence fallback.

Thus MPC shares the decision rate and command bounds with VGCC but not its model
or data budget. Its 81 sequence evaluations are numerically close to VGCC's
approximately 70 one-step candidates, but each sequence uses four recursive
ensemble transitions plus a terminal-value query. We describe it as a heavier
model-based baseline, not a same-model or equal-compute comparison.

\textbf{Mixed-terrain evaluation (Section~\ref{sec:results-headroom}).}
\label{app:hetero}
The flat and rough single-type tasks replace the terrain generator of the
otherwise unchanged rough environment with a single sub-terrain each: a plane,
and uniform height noise of $2$--$10$\,cm; every scene, observation, command,
and event setting is identical, and the same pretrained checkpoint is used.
The response model for these deployments is identified once on pooled excitation
from both terrains with the unchanged recipe ($\approx$25k windows; displacement
$R^2=0.92$, cost $0.73$). Fixed scales $\{0.60, 0.75, 0.90, 1.00\}$ and VGCC run
from identical initial states on three seeds disjoint from all tuning. Because
the world-frame final-height criterion is dominated by stop-point ground elevation
on $\pm10$\,cm noise, completion means reaching the target without early
termination. Table~\ref{tab:hetero} reports the paired frontier: COT falls
monotonically on both terrains. The budget analysis pairs one flat with one rough
leg, sweeps a shared completion-time budget over the feasible range, and compares
the best budget-feasible uniform scale against the best per-terrain scale
allocation on the convex hulls of the per-terrain (work, time) means; the
reported interval is an episode bootstrap over both legs.

\begin{table}[h]
\centering
\caption{Single-terrain-type paired evaluation (Go2, three seeds, eight targets $\times$ two trials; completion = target reached without early termination; flat $n=48$, rough $n=36$ jointly completed episodes per scale, of 48 evaluated). COT falls monotonically with the scale on both terrains.}
\label{tab:hetero}
\begin{tabular}{lcccccc}
\toprule
& \multicolumn{3}{c}{Flat} & \multicolumn{3}{c}{Rough ($2$--$10$\,cm)} \\
\cmidrule(lr){2-4}\cmidrule(lr){5-7}
Scale & COT & Work (J) & Time (s) & COT & Work (J) & Time (s) \\
\midrule
1.00 & 1.027 & 179.6 & 1.80 & 1.670 & 290.5 & 2.16 \\
0.90 & 0.964 & 166.7 & 1.99 & 1.598 & 278.5 & 2.30 \\
0.75 & 0.897 & 152.6 & 2.36 & 1.513 & 265.2 & 2.76 \\
0.60 & 0.882 & 148.8 & 2.90 & 1.447 & 247.8 & 3.31 \\
\bottomrule
\end{tabular}
\end{table}

\textbf{Exact-replay adapter necessity audit.}
The Go2 audit of Section~\ref{sec:results-headroom} uses ten source seeds, 48 query
states per seed (480 total), and constant four-macro-step prefixes from
$\{0.75,0.90,1.00\}$. For each query and scale, the simulator is reset to the
episode start, the recorded direct-control prefix is replayed to the query, the
candidate prefix is executed, and direct control completes the episode. The
terrain curriculum is frozen; all query observations reconstruct exactly
(maximum observation error zero). This start-state replay is necessary because
Isaac Lab's scene-state reset does not restore policy observation and action
history. Seeds 854, 855, 856, 858, and 859 are the original collection; seeds
860, 862, 863, 864, and 865 were added with the same protocol
(\texttt{run\_exact\_replay\_prefix\_seeds.sh}). The exploratory four-way feature
ablation table below remains on the original five-seed subset; the main text
reports the ten-seed observation-only versus observation+$F_\eta$ comparison.

\textbf{H1 second-embodiment exact replay.}
The H1 confirmation uses the same constant-scale prefix protocol
(\texttt{run\_exact\_replay\_prefix\_seeds\_h1.sh}) with
\texttt{--ff\_height\_scan\_slice 69:256}. Because trained H1 macro-transition and
task-value models were not available, collection drives episodes with
dimension-matched untrained stubs
(\texttt{h1\_macro\_transition\_h4\_stub5.pt}, \texttt{h1\_task\_value\_stub.pt}).
These stubs affect only which query states are visited; every reported
counterfactual outcome is an exact same-state replay with observation
reconstruction error zero. The H1 treatment-effect audit uses the observation-only
feature set with \texttt{--height-scan-slice 69:256} and does not consume stub
prefix predictions. Four completed source seeds (870--873; 192 queries) are
reported in the main text; additional seeds continue under the same recipe.

The same-state oracle requires candidate success no worse than the direct branch
and total remaining time no greater than $1.10$ times direct, then minimizes
total mechanical work. The treatment-effect model is a five-member, two-layer
MLP ensemble (hidden width 64) trained for 200 epochs. Inputs contain the
deployment observation, local target, candidate scale, and differences between
candidate and direct macro-response ensemble predictions (work, progress,
minimum height, maximum tilt, and maximum angular speed; means and
uncertainties). Targets are candidate work and time relative to the direct
branch and realized eligibility. Evaluation leaves one entire source seed out.
The selector accepts a candidate only when its work upper confidence estimate is
negative, its time upper estimate meets the budget, and its eligibility
probability lower estimate is at least $0.5$.

For each held-out fold, the matched randomized mixture uses the selector's exact
marginal frequencies of scales $0.75$, $0.90$, and $1.00$, but assigns them
independently of query state. Its reported metrics are the exact expectation over
the replayed potential outcomes, so no Monte Carlo sampling noise is introduced.
The selector--mixture difference isolates state-dependent allocation from the
gain due merely to issuing slower commands more often. Scripts and aggregate
artifacts are
\texttt{analyze\_exact\_replay\_oracle.py},
\texttt{train\_exact\_replay\_treatment\_effect.py},
\texttt{exact\_replay\_prefix\_oracle\_5seed.json}, and
\texttt{exact\_replay\_adapter\_necessity\_audit\_5seed.json}.

\subsection{Treatment-Effect Feature Ablation}
\label{app:feature-ablation}

We isolate which deployment-time representation contributes to the final
decision while holding the treatment-effect ensemble, five leave-one-source-seed-out
folds, confidence rule, and matched-mixture evaluation fixed. All variants
receive the current observation summary, local target, and candidate scale.
The $f_\phi$ variant adds candidate-versus-direct means and ensemble
uncertainties from the six-output one-window response model. The $F_\eta$
variant adds the corresponding four-step macro-prefix predictions used by MPC.

\begin{table}[h]
\centering
\caption{Exploratory feature ablation for the exact-replay treatment-effect selector (240 queries). ``Allocation gain'' is selector work relative to its own frequency-matched randomized mixture; negative is better. The pre-existing main result is observation+$F_\eta$.}
\label{tab:feature-ablation}
\begin{tabular}{lccccc}
\toprule
Features & Work MAE & Active & Work vs.\ direct & Allocation gain & Violations \\
\midrule
observation only & $12.60\%$ & $8.8\%$ & $-0.43\%$ & $-0.38\%$ & $0.42\%$ \\
observation + $f_\phi$ & $12.93\%$ & $13.8\%$ & $-0.37\%$ & $-0.35\%$ & $4.17\%$ \\
observation + $F_\eta$ & $\mathbf{12.22\%}$ & $12.9\%$ & $-0.54\%$ & $-0.42\%$ & $1.67\%$ \\
observation + $f_\phi+F_\eta$ & $12.33\%$ & $12.1\%$ & $\mathbf{-0.76\%}$ & $\mathbf{-0.76\%}$ & $2.92\%$ \\
\bottomrule
\end{tabular}
\end{table}

Observation-only features recover nearly the same allocation gain as the
prespecified observation+$F_\eta$ representation ($0.38\%$ versus $0.42\%$).
Adding $f_\phi$ alone does not improve work error or allocation and increases
realized constraint violations. Combining both identified models gives the
largest point estimate, $0.76\%$, but this interaction was discovered on the
same five folds, is below one percent, and incurs more violations; it is not a
nested or independently confirmed model-selection result. We therefore retain
observation+$F_\eta$ as the main prespecified audit and report the combined row
only as a hypothesis for future confirmation. The robust conclusion is that
the treatment-effect ensemble extracts most of its limited signal from the
observation itself; neither identified model provides a decisive incremental
contribution on this dataset.

\subsection{Audit Calibration Controls}
\label{app:audit-controls}

The semi-synthetic calibration changes only candidate mechanical-work outcomes;
observations, targets, success, time, posture, source-seed splits, and the
treatment-effect learner remain unchanged. For observable controls, the
preferred scale is $0.90$ when the pre-intervention local target has
$y\ge0$ and $0.75$ otherwise, and its realized work is multiplied by
$1-\gamma$. Hidden controls draw the preferred scale independently for each
query using a deterministic seed and do not expose that draw in any feature.
The matched-mixture calculation is repeated after each modification.

Confidence intervals resample the ten source-seed folds as clusters with
100{,}000 bootstrap replicates (seed 20260717). Table~\ref{tab:control-curve}
reports the ten-seed observable-signal sweep used in the main text. Decisions
use $\kappa=5\%$; G, N, and A denote \textsc{Go}, \textsc{No-Go}, and
\textsc{Abstain}.

\begin{table}[h]
\centering
\caption{Observable-signal calibration on ten source seeds. Decision columns correspond to practical allocation-value thresholds $\delta$.}
\label{tab:control-curve}
\begin{tabular}{lccccc}
\toprule
$\gamma$ & $H_{\rm alloc}$, 95\% CI & Mean $q$ & $\delta=.25\%$ & $.50\%$ & $1.00\%$ \\
\midrule
$2\%$  & $0.53\ [0.17,0.87]\%$ & $3.12\%$ & A & A & N \\
$5\%$  & $1.26\ [0.75,1.78]\%$ & $4.58\%$ & A & A & A \\
$9\%$  & $2.32\ [1.67,2.98]\%$ & $5.42\%$ & A & A & A \\
$10\%$ & $2.61\ [1.85,3.30]\%$ & $5.42\%$ & A & A & A \\
\bottomrule
\end{tabular}
\end{table}

Relative to the earlier five-seed curve, the work intervals move as expected
once $\gamma\ge5\%$, but the violation intervals are wider and cross
$\kappa=5\%$, so the joint rule abstains throughout this sweep. A five-seed
\textsc{Go} at $\gamma=9\%$ therefore does not survive the larger cluster
sample: underpowered audits can manufacture apparent detection.

To estimate specificity, we repeat the $5\%$ hidden-signal control with ten
independent mappings (control seeds 20260720--20260729) on the same ten source
seeds. The false-\textsc{Go} counts are $3/10$, $2/10$, and $0/10$ at
$\delta=0.25\%,0.50\%,1.00\%$, respectively. At $\delta=1\%$, four mappings
return \textsc{No-Go} and six \textsc{Abstain}. This stress test motivates
reporting the entire decision curve and keeping $\delta$ no smaller than the
threshold with zero false \textsc{Go} outcomes in calibration. It does not
establish a universal $1\%$ threshold: another domain requires its own
deployment value and control-based calibration.

The control experiments are generated by
\texttt{train\_exact\_replay\_treatment\_effect.py}; source-seed intervals and
decisions are produced by
\texttt{summarize\_adapter\_value\_audit.py}. Aggregate artifacts are
\texttt{adapter\_value\_audit\_observable\_curve\_10seed\_summary.json} and
\texttt{adapter\_value\_audit\_hidden\_05pct\_repeats\_10seed\_summary.json}.
The prefix collection recipe is
\texttt{run\_exact\_replay\_prefix\_seeds.sh}.

\section{Secondary Experiments and Diagnostics}
\label{app:secondary}

This appendix retains material that is useful for diagnosis and reproducibility
but is not required for the main claim: qualitative and per-target results, a
retrieval-based alternative, and the complete component ablation.

\subsection{Qualitative and Per-Target Results}
\label{app:qualitative}

\begin{figure}[h]
\centering
\includegraphics[width=\linewidth]{ff_v4_harder_trajectories.png}
\caption{XY paths on the eight harder targets (one evaluation seed, 10 trials per method). VGCC and direct control reach similar target neighborhoods.}
\label{fig:traj}
\end{figure}

\begin{figure}[h]
\centering
\includegraphics[width=\linewidth]{ff_v4_control_example.png}
\caption{Selected paired episode illustrating rescue behavior. This example shows mechanism, not prevalence.}
\label{fig:example}
\end{figure}

The selected example forecasts a posture-floor breach and activates rescue, but
the aggregate residual failure mode is low-posture arrival. A symmetric
posture-aware stopping rule did not improve success at $n=400$ per method; it
instead spent the remaining budget attempting recovery. Individual examples must
therefore not be interpreted as a success-rate effect; the aggregate evidence is
in Tables~\ref{tab:physical-energy} and~\ref{tab:coverage}.

Table~\ref{tab:pertarget} gives the per-target breakdown of the harder benchmark.
The proxy reduction holds on seven of eight targets, but success differences change
sign at $n=50$ per target, which is why we report aggregate success rather than a
target-specific success effect.

\begin{table}[h]
\centering
\caption{Per-target harder-target results (50 episodes per target and method), ordered by proxy reduction.}
\label{tab:pertarget}
\begin{tabular}{lccccccc}
\toprule
& \multicolumn{3}{c}{Effort proxy $\downarrow$} &
\multicolumn{2}{c}{Success $\uparrow$} & \multicolumn{2}{c}{Final height (m)} \\
\cmidrule(lr){2-4}\cmidrule(lr){5-6}\cmidrule(lr){7-8}
Target & VGCC & direct & $\Delta$ & VGCC & direct & VGCC & direct \\
\midrule
$(0,1.25)$ & 1.815 & 2.083 & $-12.9\%$ & 0.940 & 0.980 & 0.340 & 0.351 \\
$(1.75,-0.75)$ & 2.217 & 2.542 & $-12.8\%$ & 0.480 & 0.520 & 0.276 & 0.285 \\
$(-1,0.75)$ & 1.861 & 2.122 & $-12.3\%$ & 1.000 & 1.000 & 0.358 & 0.367 \\
$(0,-1.25)$ & 1.822 & 2.073 & $-12.1\%$ & 0.940 & 0.980 & 0.345 & 0.350 \\
$(1.5,1)$ & 2.203 & 2.471 & $-10.8\%$ & 0.800 & 0.740 & 0.316 & 0.313 \\
$(2,0)$ & 2.367 & 2.619 & $-9.6\%$ & 0.240 & 0.260 & 0.233 & 0.224 \\
$(1,-1)$ & 2.028 & 2.240 & $-9.5\%$ & 0.980 & 0.960 & 0.355 & 0.349 \\
$(-0.75,0)$ & 1.496 & 1.510 & $-0.9\%$ & 1.000 & 1.000 & 0.359 & 0.363 \\
\bottomrule
\end{tabular}
\end{table}

\subsection{Retrieval Alternative (VGSR)}
\label{app:retrieval}

An earlier, nonparametric alternative to the response model retrieves reusable
behavior chunks instead of predicting command responses. VGSR segments
converged-policy rollouts into chunks, describes them by body-frame displacement,
velocity, yaw, posture, effort proxy, stability, and smoothness, then
novelty-filters and clusters them into a behavior archive. At deployment it
retrieves a chunk by progress, state applicability, effort, stability, and utility,
and blends the chunk velocity with the direct command only when a heuristic gate
accepts it. Unlike VGCC, it can interpolate only inside archive support.

On the original five-target protocol, VGSR matches VGCC's proxy level on the harder
family (2.061 versus 2.073) but has 1.2\% terminations and lower final height. It
\emph{increases} the proxy by 3.8\% on held-out directions and 4.4\% on H1, exactly
the generalization the parametric model provides and the archive does not. In a
separate five-seed harder-target aggregate it reduces the proxy by 7.8\% with
success 0.785 versus 0.800. Removing its gate increases the proxy by more than
50\%. Archive size is non-monotonic: the top 2,000 chunks improve success relative
to all 3,904, while top-500 minimizes proxy at lower success. These coverage and
curation dependencies motivate the parametric response model used by VGCC.

\subsection{Complete VGCC Component Ablation}
\label{app:ablation}

\begin{table}[h]
\centering
\caption{Paired component ablations on two seeds ($n=80$ per variant). Differences are against direct control in the same runs.}
\label{tab:vgcc-ablation}
\begin{tabular}{lccc}
\toprule
Variant & $\Delta$ proxy $\downarrow$ & $\Delta$ success & Active steps \\
\midrule
VGCC (full) & $-12.2\%$ & $+0.012$ & 11.7 \\
no annealing & $-12.2\%$ & $+0.013$ & 12.3 \\
no viability gate & $-4.8\%$ & $-0.025$ & 8.9 \\
no progress floor & $-19.6\%$ & $-0.037$ & 19.8 \\
substitution ($\alpha=1$) & $-14.0\%$ & $-0.050$ & 16.3 \\
25\% identification data & $-6.4\%$ & $-0.012$ & 10.7 \\
\bottomrule
\end{tabular}
\end{table}

The gate keeps predicted savings inside the model's well-behaved posture regime;
the progress floor limits behavioral drift; and bounded blending avoids the larger
success cost of substitution. A quarter of the identification data retains
displacement $R^2=0.92$ but reduces the cost-channel $R^2$ from 0.75 to 0.62,
halving the realized proxy reduction. An absolute proxy margin fails across
embodiments (it increases the H1 proxy by 13.5\%), so the relative margin in
Eq.~\ref{eq:vgc-cost} is retained.

\textbf{Objective ablation.}
A variant that ranks candidates by predicted cost \emph{per predicted progress}
(with execution-aligned scoring, all other mechanisms unchanged) was evaluated
paired against the standard gate on the locked frontier protocol (three seeds,
$n=48$ episodes per method, identical initial states). It forfeits most of the
compensation: effort proxy $2.010$ versus $1.810$ for the raw-cost gate (direct
$2.132$), mechanical work $132.3$ versus $126.9$\,J, success $0.792$ versus
$0.854$, and it intervenes less (elapsed time $2.08$ versus $2.25$\,s). The
per-seed paired proxy difference favors the raw-cost gate on all three seeds.
With the progress floor already capping slowdown, dividing by predicted progress
double-penalizes slower candidates, so the variant rejects interventions the
floor would have certified; the anti-slowdown role belongs to the constraint,
not the objective.
